\documentclass[10pt]{article}
\usepackage{newspaper}
\usepackage[margin=0.55in]{geometry}
\usepackage{multicol}
\usepackage{booktabs}
\usepackage{xcolor}
\usepackage{enumitem}
\usepackage[hidelinks]{hyperref}
\SetPaperName{Neighbourhood Review}
\SetHeaderName{Neighbourhood Review}
\SetPaperLocation{Your City}
\SetPaperSlogan{``Small changes, clearly reported.''}
\SetPaperPrice{Community edition}
\currentvolume{2}
\currentissue{4}
\date{August 2026}
\setlist{nosep,leftmargin=*}
\begin{document}
\maketitle
\begin{multicols}{3}
\byline{A better route to the market}{The Field Notes Desk}
The most useful local updates are short enough to read at the corner shop and
specific enough to change what happens next. This issue follows one walking
route, three observations, and a small experiment with shade markers.

\headline{Three things we noticed}
\begin{enumerate}
  \item People paused longest where shade and a seat were separated by more than
        a few steps.
  \item A clear route marker prompted more questions than a dense map.
  \item The busiest crossing was also the place where the queue moved slowest.
\end{enumerate}

\headline{The two-week test}
Volunteers placed a movable shade marker at three waiting areas and recorded
arrival time, approximate queue length, and whether people stayed or walked on.
The measure is modest by design: it supports a next decision without claiming
to be a complete evaluation.

\begin{center}
\begin{tabular}{@{}lrr@{}}\toprule
Location & Stayed & Median wait\\\midrule
North gate & 22/31 & 6 min\\
Market corner & 19/28 & 8 min\\
Library steps & 20/24 & 4 min\\\bottomrule
\end{tabular}
\end{center}

\headline{What comes next}
Repeat the observation after rain. Bring one unanswered question and one idea
small enough to try before the following issue.

\closearticle
\headline{From the noticeboard}
\textbf{Open walk}\par
Wednesday, 18:00 — meet by the library steps for a 30-minute route review.

\medskip\textbf{Share a signal}\par
Send a short observation, a date, and a place to
\href{mailto:notes@example.org}{notes@example.org}.
\end{multicols}
\end{document}
